% Section 7: numerical illustrations.
\section{Numerical illustrations}\label{sec:numerics}

The computations below are floating-point illustrations; no proof uses them. The scripts and their stored outputs
are in the repository named in the data availability statement. Coefficients were computed either by the Lanczos
process with full reorthogonalisation in double precision on the atoms $k\le K$, or, where a tail correction is
mentioned, from the atoms $k\le K$ together with a Gauss--Jacobi rule for the continuous tail
$\int_{K+1/2}^\infty t^{-s}g(t^{-p})\diff t$, by an algorithm that adds one atom at a time and restores tridiagonal form by
plane rotations, in quadruple precision; the latter reproduces the Lommel coefficients of
Proposition~\ref{prop:imp-lommel} to a relative error below $10^{-14}$ for $n\le400$.

\emph{Rough weights.} Table~\ref{tab:rough} shows the construction of Theorem~\ref{thm:rough-b} for $p=s=2$, the window
$W=(0.6,1.6)$, three scales $N_i=15,60,240$ (separation by the factor $4$, far below $R_W$) and two values of
$\varepsilon$, both far above the proved range $\varepsilon\le\gamma_W/16\approx6\cdot10^{-9}$. At each scale the jump
$F_m(w)-F_m(v)$ at the selected index $m$ agrees with the first-order prediction $\langle\varphi,\Delta^2m_m^v\rangle$
to within $2\%$, the contribution of the other windows is at most $1.4\cdot10^{-3}$, and every jump exceeds the
largest $|F_n(w)-F_n(v)|$ over the indices $n\in[10,246]$ with $|n-m_i|>2$ for all $i$ by a factor larger than $2$. The measure was truncated at
$K=9600$ atoms; with two scales $(30,240)$ and $\varepsilon=0.2$, doubling $K$ changes the jumps by less than
$10^{-3}$.

\begin{table}[ht]
\centering
\small
\begin{tabular}{rrrrrr}
\toprule
$\varepsilon$ & $N_i$ & $m_i$ & $F_{m_i}(w)-F_{m_i}(v)$ & $\langle\varphi,\Delta^2m^v_{m_i}\rangle$ & other windows\\
\midrule
0.05 & 15  & 15  & $0.0689$ & $0.0691$ & $-3.5\cdot10^{-4}$\\
0.05 & 60  & 58  & $0.0624$ & $0.0626$ & $2.0\cdot10^{-4}$\\
0.05 & 240 & 238 & $0.0601$ & $0.0601$ & $<10^{-16}$\\
0.2  & 15  & 15  & $0.2739$ & $0.2765$ & $-1.4\cdot10^{-3}$\\
0.2  & 60  & 58  & $0.2457$ & $0.2503$ & $8.2\cdot10^{-4}$\\
0.2  & 240 & 238 & $0.2389$ & $0.2403$ & $<10^{-16}$\\
\bottomrule
\end{tabular}
\caption{The construction of Theorem~\ref{thm:rough-b} for $p=s=2$, $W=(0.6,1.6)$, scales $15,60,240$. The largest
$|F_n(w)-F_n(v)|$ over $n\in[10,246]$ with $|n-m_i|>2$ is $0.026$ for $\varepsilon=0.05$ and $0.103$ for
$\varepsilon=0.2$. Double precision, $K=9600$ atoms.}
\label{tab:rough}
\end{table}

\emph{The Christoffel--Ces\`aro identity and the $s$-law.} Table~\ref{tab:shift} lists the error
$S_n:=\sum_{m\le n}\log(b_m(s+p)/b_m(s))+\frac p2\log n-\frac12\log(c\zeta(s)/\zeta(s+p))$ in Theorem~\ref{thm:s}(i),
computed with the tail correction ($K=3.2\cdot10^5$ atoms, and $K=6.4\cdot10^5$ for $p=\frac12$). It decays like $1/n$, much faster than the
a priori rate $n^{-1/2}$. For the uniform statements, on the grid $s\in\{1.5,1.75,\dots,4\}$ and with $K=4\cdot10^4$ plus the
tail correction, the quantity $\sup_s|\Phi_n(s)-\Phi_n(\frac32)-(s-\frac32)\log c/p|$ of Theorem~\ref{thm:UFlim} equals
$2.25\cdot10^{-2}$, $1.12\cdot10^{-2}$, $5.62\cdot10^{-3}$, $2.81\cdot10^{-3}$ for $p=1$ and $5.26\cdot10^{-2}$,
$2.64\cdot10^{-2}$, $1.32\cdot10^{-2}$, $6.63\cdot10^{-3}$ for $p=2$, at $n=50,100,200,400$; the deviations
$\sup_s|n^pb_n(s)/c-1|$ and $\sup_s|n^pa_n(s)/(2c)-1|$ are $3.7\cdot10^{-3}$ and $5.0\cdot10^{-3}$ ($p=1$), and
$3.7\cdot10^{-3}$ and $6.2\cdot10^{-3}$ ($p=2$), at $n=400$, and halve with each doubling of $n$. For $p=1$ the grid extends
beyond $\Sp=(1,3)$; there the computation illustrates what we expect but have not proved. Finally, for the
Tricomi--Carlitz comparison in the proof of Corollary~\ref{cor:lattice}(a), the error of
$\sum_{j\le n}\log(b_j(\frac32)/b^{\mathrm{TC}}_j)$, computed with the value $W_{\mathrm{TC}}=1$ that the computation
indicates, is $4.1\cdot10^{-4}$ at $n=100$ and $2.6\cdot10^{-5}$ at $n=1600$.

\begin{table}[ht]
\centering
\small
\begin{tabular}{rrrrrr}
\toprule
$p$ & $s$ & $n=100$ & $n=400$ & $n=800$ & $n=1599$\\
\midrule
0.5 & 1.5 & $-6.4\cdot10^{-5}$ & $-6.7\cdot10^{-6}$ & $-2.2\cdot10^{-6}$ & $-7.6\cdot10^{-7}$\\
1   & 2   & $-2.5\cdot10^{-3}$ & $-6.3\cdot10^{-4}$ & $-3.1\cdot10^{-4}$ & $-1.6\cdot10^{-4}$\\
1   & 3   & $2.3\cdot10^{-3}$  & $6.1\cdot10^{-4}$  & $3.1\cdot10^{-4}$  & $1.6\cdot10^{-4}$\\
1.5 & 1.5 & $-8.3\cdot10^{-3}$ & $-2.1\cdot10^{-3}$ & $-1.0\cdot10^{-3}$ & $-5.2\cdot10^{-4}$\\
2   & 2   & $-9.9\cdot10^{-3}$ & $-2.5\cdot10^{-3}$ & $-1.2\cdot10^{-3}$ & $-6.3\cdot10^{-4}$\\
\bottomrule
\end{tabular}
\caption{The error $S_n$ in the Christoffel--Ces\`aro identity, Theorem~\ref{thm:s}(i). Quadruple precision, atoms
plus a quadrature of the continuous tail.}
\label{tab:shift}
\end{table}

\emph{Slowly varying factors.} For $p=s=2$ and the unbounded factor $\ell(k)=\exp\sqrt{\log(k+e)}$, the relative
deviation $n^pb_n(\sigma_{p,s,\ell})/(n^pb_n(\sigma_{p,s}))-1$ equals $1.09\cdot10^{-3}$, $5.23\cdot10^{-4}$,
$2.49\cdot10^{-4}$, $1.18\cdot10^{-4}$, $5.6\cdot10^{-5}$ at $n=100,200,400,800,1600$ (extended precision, $K=1.2\cdot10^5$
and $2.4\cdot10^5$ atoms without tail correction, which agree to $2\cdot10^{-7}$); $n$ times it decreases slowly, like
$\frac12\frac{\diff}{\diff\xi}\sqrt{\log(e^\xi+e)}$ at $\xi=\log n$. For the factor $\ell(k)=1+\frac12(-1)^k/\log(k+1)$ of
Remark~\ref{rem:rate}, which is slowly varying but rough at every scale, the deviation $|b_n(\sigma_{2,2,\ell})/b_n(\sigma_{2,2})-1|$
is $3.9\cdot10^{-3}$, $1.5\cdot10^{-3}$ and $8.2\cdot10^{-4}$ at $n=250,1000,3600$ (extended precision, a single truncation
at $K=1.2\cdot10^5$ atoms, so these values are only indicative); the decay is slow, as the rate $\rho_\ell(\lceil
n/M_0\rceil)=O(1/\log n)$ of Theorem~\ref{thm:sv} allows.
